\begin{tomtat}
	\textbf{ĐỀ 02}
\end{tomtat}
\begin{ex}
	Rút gọn biểu thức $A = \left( \dfrac{\sqrt{x}-1}{\sqrt{x}+1} - \dfrac{\sqrt{x}+1}{\sqrt{x}-1} - \dfrac{4x}{x-1} \right) . \dfrac{2(\sqrt{x}-1)}{\sqrt{x}+4}$ (với $x \ge 0, x \ne 1$).
	\loigiai{
		Với điều kiện $x \ge 0, x \ne 1$, ta có:
		\begin{align*}
			A &= \left( \dfrac{\sqrt{x}-1}{\sqrt{x}+1} - \dfrac{\sqrt{x}+1}{\sqrt{x}-1} - \dfrac{4x}{x-1} \right) . \dfrac{2(\sqrt{x}-1)}{\sqrt{x}+4} \\
			&= \left( \dfrac{(\sqrt{x}-1)^2}{(\sqrt{x}+1)(\sqrt{x}-1)} - \dfrac{(\sqrt{x}+1)^2}{(\sqrt{x}-1)(\sqrt{x}+1)} - \dfrac{4x}{(\sqrt{x}-1)(\sqrt{x}+1)} \right) \cdot \dfrac{2(\sqrt{x}-1)}{\sqrt{x}+4} \\
			&= \dfrac{(x - 2\sqrt{x} + 1) - (x + 2\sqrt{x} + 1) - 4x}{(\sqrt{x}-1)(\sqrt{x}+1)} \cdot \dfrac{2(\sqrt{x}-1)}{\sqrt{x}+4} \\
			&= \dfrac{x - 2\sqrt{x} + 1 - x - 2\sqrt{x} - 1 - 4x}{x-1} \cdot \dfrac{2(\sqrt{x}-1)}{\sqrt{x}+4} \\
			&= \dfrac{-4\sqrt{x} - 4x}{x-1} \cdot \dfrac{2(\sqrt{x}-1)}{\sqrt{x}+4} \\
			&= \dfrac{-4\sqrt{x}(1+\sqrt{x})}{(\sqrt{x}-1)(\sqrt{x}+1)} \cdot \dfrac{2(\sqrt{x}-1)}{\sqrt{x}+4} \\
			&= \dfrac{-4\sqrt{x}}{\sqrt{x}-1} \cdot \dfrac{2(\sqrt{x}-1)}{\sqrt{x}+4} \\
			&= \dfrac{-4\sqrt{x} \cdot 2(\sqrt{x}-1)}{(\sqrt{x}-1)(\sqrt{x}+4)} \\
			&= \dfrac{-8\sqrt{x}}{\sqrt{x}+4}
		\end{align*}
		Vậy biểu thức rút gọn là $A = \dfrac{-8\sqrt{x}}{\sqrt{x}+4}$.
	}
\end{ex}


\begin{ex}
	Giải hệ phương trình $\begin{cases} 2x+3y=-1 \\ 2x-y=-5 \end{cases}$.
	\loigiai{
		Ta có hệ phương trình:
		$$ \begin{cases} 2x+3y=-1 & (1) \\ 2x-y=-5 & (2) \end{cases} $$
		Lấy phương trình (1) trừ đi phương trình (2) vế theo vế, ta được:
		\begin{align*}
			(2x+3y) - (2x-y) &= -1 - (-5) \\
			2x+3y-2x+y &= -1+5 \\
			4y &= 4 \\
			y &= 1
		\end{align*}
		Thay $y=1$ vào phương trình (2), ta được:
		\begin{align*}
			2x - 1 &= -5 \\
			2x &= -5+1 \\
			2x &= -4 \\
			x &= -2
		\end{align*}
		Vậy nghiệm của hệ phương trình là $(x; y) = (-2; 1)$.
	}
\end{ex}


\begin{ex}
	Cho phương trình $x^2 - 2x - 1 = 0$, không giải phương trình, hãy tính giá trị biểu thức $T = (x_1+x_2)^2 - x_1^2x_2^2$.
	\loigiai{
		Phương trình đã cho là $x^2 - 2x - 1 = 0$.
		Đây là phương trình bậc hai có dạng $ax^2+bx+c=0$ với $a=1$, $b=-2$, $c=-1$.
		Ta tính biệt thức $\Delta = b^2-4ac = (-2)^2 - 4(1)(-1) = 4 + 4 = 8$.
		Vì $\Delta = 8 > 0$ nên phương trình có hai nghiệm phân biệt $x_1, x_2$.
		
		Theo hệ thức Vi-ét, ta có:
		$$ \begin{cases} 
			x_1+x_2 = \dfrac{-b}{a} = \dfrac{-(-2)}{1} = 2 \\ 
			x_1x_2 = \dfrac{c}{a} = \dfrac{-1}{1} = -1 
		\end{cases} $$
		Ta cần tính giá trị biểu thức $T = (x_1+x_2)^2 - x_1^2x_2^2$.
		Biểu thức $T$ có thể viết lại dưới dạng:
		$T = (x_1+x_2)^2 - (x_1x_2)^2$.
		Thay giá trị $x_1+x_2$ và $x_1x_2$ đã tính ở trên vào biểu thức $T$:
		\begin{align*}
			T &= (2)^2 - (-1)^2 \\
			&= 4 - 1 \\
			&= 3
		\end{align*}
		Vậy giá trị của biểu thức $T$ là 3.
	}
\end{ex}


\begin{ex} % Câu 13 (Xúc xắc)
	Gieo 1 con xúc xắc cân đối đồng chất một lần. Tính xác suất của biến cố A: "mặt xuất hiện có số chấm là số nguyên tố".
	\loigiai{
		Tập hợp tất cả các kết quả có thể xảy ra khi gieo một con xúc xắc (không gian mẫu $\Omega$) là:
		$\Omega = \{1, 2, 3, 4, 5, 6\}$.
		Số phần tử của không gian mẫu là $n(\Omega) = 6$.
		
		Tập hợp các kết quả thuận lợi cho biến cố A: "mặt xuất hiện có số chấm là số nguyên tố".
		Các số nguyên tố trong các mặt của xúc xắc là: 2, 3, 5.
		Vậy, $A = \{2, 3, 5\}$.
		Số phần tử của biến cố A là $n(A) = 3$.
		
		Xác suất của biến cố A là:
		$$ P(A) = \frac{n(A)}{n(\Omega)} = \frac{3}{6} = \frac{1}{2} $$
		Vậy xác suất của biến cố A là $\frac{1}{2}$.
	}
\end{ex}
\begin{ex}
	Bà Hương vay 500 triệu đồng của ngân hàng trong thời hạn 3 năm để mở một cửa hàng bán đồ thủ công mỹ nghệ. Theo hợp đồng vay vốn, lãi suất vay trong 1 năm là 7\%. Sau mỗi năm, tiền lãi của năm trước sẽ được cộng vào vốn vay của năm sau. Sau 2 năm, bà Hương phải trả cho ngân hàng số tiền cả gốc lẫn lãi là bao nhiêu?
	\loigiai{
		Số tiền vốn ban đầu bà Hương vay là: $P_0 = 500,000,000$ đồng.
		Lãi suất vay mỗi năm là: $r = 7\% = 0.07$.
		Bài toán này áp dụng công thức lãi kép, theo đó tiền lãi của kỳ trước được cộng vào vốn để tính lãi cho kỳ tiếp theo.
		
		Số tiền cả gốc lẫn lãi bà Hương phải trả sau năm thứ nhất là:
		\begin{align*}
			P_1 &= P_0(1+r) \\
			&= 500,000,000 \times (1 + 0.07) \\
			&= 500,000,000 \times 1.07 \\
			&= 535,000,000 \text{ đồng}
		\end{align*}
		
		Số tiền cả gốc lẫn lãi bà Hương phải trả sau năm thứ hai là (với vốn đầu năm thứ hai là $P_1$):
		\begin{align*}
			P_2 &= P_1(1+r) \\
			&= 535,000,000 \times (1 + 0.07) \\
			&= 535,000,000 \times 1.07 \\
			&= 572,450,000 \text{ đồng}
		\end{align*}
		
		Cách khác, ta có thể tính trực tiếp số tiền sau 2 năm bằng công thức lãi kép:
		$P_n = P_0(1+r)^n$, với $n=2$.
		\begin{align*}
			P_2 &= 500,000,000 \times (1 + 0.07)^2 \\
			&= 500,000,000 \times (1.07)^2 \\
			&= 500,000,000 \times 1.1449 \\
			&= 572,450,000 \text{ đồng}
		\end{align*}
		Vậy, sau 2 năm, bà Hương phải trả cho ngân hàng số tiền cả gốc lẫn lãi là 572,450,000 đồng.
	}
\end{ex}


\begin{ex}
	Cho đường tròn $(O)$, bán kính $R (R>0)$ và dây cung $BC$ cố định. Một điểm $A$ chuyển động trên cung lớn $BC$ sao cho tam giác $ABC$ có ba góc nhọn. Kẻ các đường cao $AD, BE$ của tam giác $ABC$ cắt nhau tại $H$ và $BE$ cắt đường tròn $(O)$ tại $F$ ($F \neq B$).
	\begin{enumerate}
		\item Chứng minh rằng tứ giác $DHEC$ nội tiếp.
		%\item Kẻ đường kính $AM$ của đường tròn $(O)$ và $OI$ vuông góc với $BC$ tại $I$. Chứng minh rằng $I$ là trung điểm của $HM$.
		%\item Khi $BC$ cố định, xác định vị trí của $A$ trên đường tròn $(O)$ để $DH \cdot DA$ lớn nhất.
	\end{enumerate}
	\loigiai{
		\begin{center}
			\includegraphics[scale=1]{hinh/2}	
		\end{center}
			\begin{enumerate}
			\item Chứng minh rằng tứ giác $DHEC$ nội tiếp.
			\item Kẻ đường kính $AM$ của đường tròn $(O)$ và $OI$ vuông góc với $BC$ tại $I$. Chứng minh rằng $I$ là trung điểm của $HM$.
			
				
			APHI là tứ giác nội tiếp (câu 1) nên $\widehat{IAH} = \widehat{IPH}$.
			
			Tương tự ta cũng có BQHI là tứ giác nội tiếp nên:
			$\widehat{IBH} = \widehat{IQH}$ (các góc nội tiếp cùng chắn cung IH)
			
			Do đó, $\triangle AHB \sim \triangle PIQ$ nên $\widehat{AHB} = \widehat{PIQ}$. Mà $\widehat{PIQ} = 90^\circ$ nên $\widehat{AHB} = 90^\circ$.
			$H \in (O)$ và MHNI là tứ giác nội tiếp nên $\widehat{HNM} = \widehat{HIM} = \widehat{HAP}$
			
			Vì $\widehat{HAP} = \widehat{HBA}$ nên $\widehat{HNM} = \widehat{HBA}$.
			
			Mà 2 góc này ở vị trí đồng vị nên $MN // AB$.
			\item Khi $BC$ cố định, xác định vị trí của $A$ trên đường tròn $(O)$ để $DH \cdot DA$ lớn nhất.
			
					$\triangle API \sim \triangle BIQ$ suy ra: $\frac{AP}{BI} = \frac{AI}{BQ} \Rightarrow AP \cdot BQ = AI \cdot BI$ không đổi.
			
			$S_{\triangle PIQ} = \frac{1}{2}IP \cdot IQ = \frac{1}{2}\sqrt{AP^2 + AI^2} \cdot \sqrt{BQ^2 + BI^2}$
			
			$\ge \frac{1}{2}\sqrt{2AP \cdot AI} \cdot \sqrt{2BQ \cdot BI} = \sqrt{AP \cdot BQ \cdot AI \cdot BI} = AI \cdot BI = \frac{R}{2} \cdot \frac{3R}{2} = \frac{3R^2}{4}$
			
			Dấu “=” xảy ra khi $AP = AI, BQ = BI$. Vậy $S_{\triangle IPQ}$ đạt GTNN bằng $\frac{3R^2}{4}$.
		\end{enumerate}
	
		}
\end{ex}
